% Generated by scripts/build_latex.py from docs/report/ch2_literature_review.md: edit the Markdown, not this file.
\chapter{Literature Review}\label{ch:2}

This chapter reviews the work on which IDSS builds: the evolution of decision support systems, automated machine learning, LLM agents, retrieval-augmented generation, explainable AI and counterfactual explanations. It then describes AutoML-Agent, the framework that inspired IDSS, compares related work and identifies the research gap.

\FloatBarrier
\section{Evolution of Decision Support Systems}\label{sec:2.1}

Power defines a decision support system as an interactive computer-based system that helps decision makers use communication technologies, data, documents, knowledge and models to identify and solve problems \cite{ref1}. He groups such systems by their dominant component into \textbf{data-driven}, \textbf{model-driven}, \textbf{knowledge-driven}, \textbf{document-driven} and \textbf{communication-driven} systems \cite{ref1}. Table~\ref{tab:2.1} summarises, in broad terms, how decision support has evolved. Each stage made more data and stronger models available, but each still kept an expert in the loop to build the model, write the query or configure the dashboard.

\begingroup
\linespread{1}\small
\setlength{\tabcolsep}{4pt}
\begin{longtable}{|>{\raggedright\arraybackslash\hspace{0pt}}p{\dimexpr 0.2425\linewidth-2\tabcolsep-2\arrayrulewidth\relax}|>{\raggedright\arraybackslash\hspace{0pt}}p{\dimexpr 0.2619\linewidth-2\tabcolsep-2\arrayrulewidth\relax}|>{\raggedright\arraybackslash\hspace{0pt}}p{\dimexpr 0.2870\linewidth-2\tabcolsep-2\arrayrulewidth\relax}|>{\raggedright\arraybackslash\hspace{0pt}}p{\dimexpr 0.2086\linewidth-2\tabcolsep-2\arrayrulewidth\relax}|}
\caption{Broad stages in the evolution of decision support}\label{tab:2.1}\\
\hline
\rowcolor{tableheader}\textbf{Stage} & \textbf{Typical tools} & \textbf{What the user receives} & \textbf{Who builds the analysis} \\
\hline
\endfirsthead
\hline
\rowcolor{tableheader}\textbf{Stage} & \textbf{Typical tools} & \textbf{What the user receives} & \textbf{Who builds the analysis} \\
\endhead
Model-driven DSS & Spreadsheet, optimisation and simulation models & Answers from one hand-built model & Analyst \\ \hline
Data-driven DSS and BI & Data warehouses, OLAP, reports, dashboards & Descriptions of what has happened & BI developers \\ \hline
Knowledge-driven DSS & Rule-based expert systems & Rule-based recommendations & Knowledge engineers \\ \hline
Predictive analytics & Machine-learning models \cite{ref2}, \cite{ref3}, \cite{ref4} & Predictions with accuracy figures & Data scientists \\ \hline
Automated ML & AutoML systems \cite{ref5} & A tuned model for a specified task & Users who can specify target, metric and data \\ \hline
LLM-assisted DSS (this project) & LLM agents + deterministic AutoML + explanations & Verified model, explanations, actions and a report from a plain-language question & The system, with the user answering questions \\ \hline
\end{longtable}
\endgroup

IDSS belongs to the last stage. It is model-driven at its core --- every recommendation comes from a trained and verified model --- but it uses an LLM for the parts that used to need an expert to talk to the user: understanding the question, asking for missing information and explaining the result.

\FloatBarrier
\section{Automated Machine Learning and Hyperparameter Optimisation}\label{sec:2.2}

\textbf{Auto-sklearn} framed AutoML as combined algorithm selection and hyperparameter optimisation over scikit-learn pipelines, using Bayesian optimisation, meta-learning to warm-start the search and an ensemble of the models found \cite{ref5}. It matches or beats manual model building on many tabular datasets, but the user must still specify the features, target, task and metric in code, and the output is a model, not a decision.

\textbf{TPE} models the hyperparameter densities of good and bad trials separately and proposes the configuration that maximises their ratio; it handles conditional, mixed search spaces cheaply \cite{ref6}. \textbf{Optuna} offers a \textit{define-by-run} API, efficient samplers such as TPE and pruners that stop unpromising trials early \cite{ref7}. IDSS uses Optuna's TPE sampler with a median pruner.

The IDSS model zoo is built from \textbf{scikit-learn}, with its consistent \code{fit} / \code{predict} interface and pipelines \cite{ref2}, \textbf{XGBoost}, which made gradient tree boosting scalable through a regularised objective and sparsity-aware split finding \cite{ref3}, and \textbf{LightGBM}, which made it faster with histogram-based, leaf-wise growth, gradient-based one-side sampling and exclusive feature bundling \cite{ref4}.

\textbf{Gap:} AutoML automates the \textit{search}, not the \textit{framing} of the problem or the \textit{use} of the result.

\FloatBarrier
\section{Large Language Model Agents}\label{sec:2.3}

An LLM agent uses a language model to decide the next step, calls tools and reads their results. \textbf{ReAct} interleaves reasoning traces with actions, so the model can plan, act, observe and correct itself, which reduces hallucination compared with reasoning alone \cite{ref8}. \textbf{AutoGen} lets several conversable agents cooperate through messages, with a human in the loop and execution of the code they write \cite{ref12}. \textbf{Data Interpreter} decomposes data-science tasks into a graph of sub-tasks and writes and executes code for each step, revising the plan when execution fails \cite{ref13}.

These systems show that LLMs can drive multi-step analyses, but also the weakness that matters most for decision support: the reported numbers come from code written during the run, which may be wrong in ways a non-expert cannot detect. IDSS therefore lets agents \textit{choose and explain} but not \textit{compute}: they pick models only from a fixed, tested zoo, and the one place where LLM-written code runs (\textit{Ask Your Data}) is isolated in a sandbox.

\FloatBarrier
\section{Retrieval-Augmented Generation}\label{sec:2.4}

Retrieval-augmented generation (RAG) combines a generator with a non-parametric memory: a retriever fetches relevant passages from an indexed corpus and the generator conditions on them, which improves factual accuracy and lets knowledge be updated without retraining \cite{ref14}. IDSS uses a lightweight, offline form of RAG for \textit{planning}: ten best-practice notes written by the team are split into chunks of about 500 characters, indexed with TF-IDF and cached; the Knowledge Agent retrieves the four most relevant passages (at most two per note) and condenses them into a brief for the Model Agent. Web search is used only when a key is configured, so retrieval stays deterministic, fast and auditable.

\FloatBarrier
\section{Explainable AI}\label{sec:2.5}

\textbf{LIME} explains one prediction by fitting a simple model to the black-box model's behaviour on perturbed samples around that record \cite{ref15}. \textbf{SHAP} unifies additive feature attributions under Shapley values and is the only additive attribution with local accuracy, missingness and consistency \cite{ref10}. A SHAP value \textit{\ensuremath{\varphi}}\textit{\textsubscript{j}} is the contribution of feature \textit{j} to moving one prediction away from the average prediction; the mean absolute values give a global ranking. IDSS uses SHAP's tree explainer for tree ensembles and the linear explainer for logistic and ridge regression (permutation importance otherwise), sums the values of derived features back to the user's original columns --- \code{contract\_\allowbreak{}type} instead of \code{cat\_\allowbreak{}\_\allowbreak{}contract\_\allowbreak{}type\_\allowbreak{}Month-\allowbreak{}to-\allowbreak{}month} --- and orients binary explanations to the positive class, so a positive value always means ``raises the risk''.

\FloatBarrier
\section{Counterfactual Explanations and What-if Analysis}\label{sec:2.6}

Wachter, Mittelstadt and Russell proposed counterfactual explanations: the smallest change to the input that would change the model's decision \cite{ref11}. Unlike feature importances, counterfactuals directly suggest an action. IDSS implements the idea twice: a \textbf{what-if simulator} that changes inputs for one record or a segment and shows the prediction before and after, and a \textbf{counterfactual search} that finds at most three changes to actionable columns, using values observed in the data, that bring a high-risk record below the decision threshold. As in the literature, these statements describe the \textit{model}, not causal effects in the real world.

\FloatBarrier
\section{AutoML-Agent and What IDSS Changes}\label{sec:2.7}

\textbf{AutoML-Agent} \cite{ref9} is a multi-agent LLM framework for full-pipeline AutoML. An agent manager receives the request, a prompt agent turns it into a structured specification, retrieval-augmented planning produces candidate plans, data and model agents carry them out, an operation agent writes and runs the code, and verification is applied to the request, the execution and the implementation. The framework is evaluated with a success rate (SR), a normalised performance score (NPS) and a comprehensive score (CS = 0.5 \ensuremath{\cdot} SR + 0.5 \ensuremath{\cdot} NPS). We read the reference implementation for ideas only; no code was copied. Tables~\ref{tab:2.2} and~\ref{tab:2.3} list what we kept and what we changed.

\begingroup
\linespread{1}\small
\setlength{\tabcolsep}{4pt}
\begin{longtable}{|>{\raggedright\arraybackslash\hspace{0pt}}p{\dimexpr 0.2708\linewidth-2\tabcolsep-2\arrayrulewidth\relax}|>{\raggedright\arraybackslash\hspace{0pt}}p{\dimexpr 0.3393\linewidth-2\tabcolsep-2\arrayrulewidth\relax}|>{\raggedright\arraybackslash\hspace{0pt}}p{\dimexpr 0.3899\linewidth-2\tabcolsep-2\arrayrulewidth\relax}|}
\caption{Ideas kept from AutoML-Agent}\label{tab:2.2}\\
\hline
\rowcolor{tableheader}\textbf{Idea} & \textbf{In AutoML-Agent} & \textbf{In IDSS} \\
\hline
\endfirsthead
\hline
\rowcolor{tableheader}\textbf{Idea} & \textbf{In AutoML-Agent} & \textbf{In IDSS} \\
\endhead
Multi-agent decomposition & Manager with prompt, data, model and operation agents & Orchestrator with Requirement, Data, Knowledge, Model, Verifier, Analysis and Insight agents, each with typed (Pydantic) inputs and outputs \\ \hline
Structured requirement parsing & Prompt agent produces a JSON specification & Requirement Agent produces a \code{UserRequirement} and repairs it against the data \\ \hline
Retrieval-augmented planning & Plans informed by retrieved knowledge & Model Agent ranks candidates with a brief retrieved from best-practice notes \\ \hline
Multi-stage verification & Request, execution and implementation verified & Relevance and sufficiency checks; audited single use of the test split; seven result checks with one re-plan \\ \hline
Evaluation metrics & SR, NPS, CS & Reproduced offline; adjusted Rand index instead of Rand index for clustering (Section~\ref{sec:7.7}) \\ \hline
\end{longtable}
\endgroup

\begin{table}[H]
\centering
\linespread{1}\small
\setlength{\tabcolsep}{4pt}
\caption{Design decisions changed in IDSS}\label{tab:2.3}
\begin{tabular}{|>{\raggedright\arraybackslash\hspace{0pt}}p{\dimexpr 0.5043\linewidth-2\tabcolsep-2\arrayrulewidth\relax}|>{\raggedright\arraybackslash\hspace{0pt}}p{\dimexpr 0.4957\linewidth-2\tabcolsep-2\arrayrulewidth\relax}|}
\hline
\rowcolor{tableheader}\textbf{Change} & \textbf{Reason} \\
\hline
The LLM only ranks names from a fixed model zoo; tested code trains, tunes, selects and evaluates & Every run that starts can finish; results are reproducible; the test split provably stays untouched; works without an LLM \\ \hline
Offline RAG over ten curated notes (TF-IDF, cached); web search only with a key & No network dependency; no data leaves the machine; auditable retrieval \\ \hline
Sandbox for the only LLM-written code (\textit{Ask Your Data}) & LLM-generated code is untrusted input \\ \hline
Decision layer: drivers on original columns, what-if, risk tiers, counterfactual actions, grounded brief, PDF report & Decision makers need to know \textit{what to do}, not only how accurate a model is \\ \hline
Persistent, resumable runs with a REST API and a web UI & Non-experts use a browser; runs survive restarts and can wait for answers \\ \hline
Fast / full budgets and parsed constraints (latency, training time, explainability) & Predictable runtime on a laptop; constraints are enforced, not only mentioned \\ \hline
\end{tabular}
\end{table}

\FloatBarrier
\section{Comparison of Related Work}\label{sec:2.8}

\begingroup
\linespread{1}\small
\setlength{\tabcolsep}{4pt}
\begin{longtable}{|>{\raggedright\arraybackslash\hspace{0pt}}p{\dimexpr 0.1870\linewidth-2\tabcolsep-2\arrayrulewidth\relax}|>{\raggedright\arraybackslash\hspace{0pt}}p{\dimexpr 0.2130\linewidth-2\tabcolsep-2\arrayrulewidth\relax}|>{\raggedright\arraybackslash\hspace{0pt}}p{\dimexpr 0.2130\linewidth-2\tabcolsep-2\arrayrulewidth\relax}|>{\raggedright\arraybackslash\hspace{0pt}}p{\dimexpr 0.1739\linewidth-2\tabcolsep-2\arrayrulewidth\relax}|>{\raggedright\arraybackslash\hspace{0pt}}p{\dimexpr 0.2130\linewidth-2\tabcolsep-2\arrayrulewidth\relax}|}
\caption{Comparison of related work with IDSS}\label{tab:2.4}\\
\hline
\rowcolor{tableheader}\textbf{Work} & \textbf{Approach} & \textbf{Strengths} & \textbf{Limitations} & \textbf{How IDSS differs} \\
\hline
\endfirsthead
\hline
\rowcolor{tableheader}\textbf{Work} & \textbf{Approach} & \textbf{Strengths} & \textbf{Limitations} & \textbf{How IDSS differs} \\
\endhead
Model- and data-driven DSS \cite{ref1} & Hand-built models, reports, dashboards & Trusted, domain-specific & Each new question needs an analyst; mostly descriptive & Builds a predictive model per plain-language question \\ \hline
Auto-sklearn \cite{ref5} & Bayesian optimisation over pipelines, meta-learning, ensembles & Strong automatic search on tabular data & Problem specified in code; output is a model & Parses the question; adds verification and decisions \\ \hline
TPE / Optuna \cite{ref6}, \cite{ref7} & Model-based hyperparameter search with pruning & Efficient, flexible & A library: objective and models must be written & Used as the tuning engine of a tested model zoo \\ \hline
ReAct \cite{ref8} & Interleaved reasoning and tool use & General; grounds reasoning in tool results & No notion of splits, baselines or model quality & Fixed state machine with bounded loops; tools are tested engines \\ \hline
AutoGen \cite{ref12} & Conversing agents that run code & Flexible workflows, human in the loop & Results depend on generated code; hard to bound & Typed hand-offs, \ensuremath{\leq} 2 clarification rounds, 1 re-plan \\ \hline
Data Interpreter \cite{ref13} & Task graph with code per step & End-to-end data analyses & Numbers come from code generated during the run & Numbers only from deterministic code; LLM code sandboxed \\ \hline
AutoML-Agent \cite{ref9} & Multi-agent LLM framework, request to deployable model & Full pipeline, systematic verification & Generated training code can fail; stops at the model & Deterministic engine, offline fallbacks, decision layer, UI \\ \hline
SHAP / LIME \cite{ref10}, \cite{ref15} & Post-hoc feature attributions & Well-founded explanations & Expert libraries; engineered feature names & Built in, on original columns, in plain language \\ \hline
\end{longtable}
\endgroup

\FloatBarrier
\section{Research Gap}\label{sec:2.9}

AutoML systems automate model search but need a precisely framed problem and return only a model. LLM agents understand plain language but compute their numbers with code generated during the run. Explanation and counterfactual methods exist but are rarely combined with model building in a tool for non-experts. We found few systems that (1) accept a plain-language question and repair it against the data, (2) compute every number with deterministic, tested code and work without an LLM, (3) audit their own evaluation and say when a result is not verified, (4) turn the model into decisions with a brief whose numbers are checked against computed facts, and (5) run as a resumable web application on a single laptop. IDSS is designed to fill this gap; its contribution is the combination and the division of labour, not a new learning algorithm.

\FloatBarrier
\section{Summary}\label{sec:2.10}

This chapter traced decision support from hand-built models to AutoML and LLM agents and reviewed RAG, SHAP, LIME and counterfactual explanations. We described AutoML-Agent, listed the ideas we kept and the decisions we changed, compared related work and identified the gap that IDSS addresses: a decision support system that is conversational at the front, deterministic in its computation, self-verifying and decision-oriented.
